\documentclass[11pt]{article}
\usepackage{amsmath,amsthm,amssymb}
\usepackage[margin=1in]{geometry}
\usepackage{hyperref}
\usepackage{booktabs}
\usepackage{enumitem}
\usepackage{amssymb}

\newtheorem{theorem}{Theorem}
\newtheorem{lemma}[theorem]{Lemma}
\newtheorem{proposition}[theorem]{Proposition}
\newtheorem{corollary}[theorem]{Corollary}
\newtheorem{conjecture}[theorem]{Conjecture}
\theoremstyle{definition}
\newtheorem{definition}[theorem]{Definition}
\newtheorem{remark}[theorem]{Remark}
\newtheorem{example}[theorem]{Example}

\newcommand{\Hq}{H_q}
\newcommand{\Sn}{S_n}
\newcommand{\Vlam}{V_\lambda}
\newcommand{\Bdec}{B_{\ell+1}^{(\lambda)}}
\newcommand{\Om}{\Omega}
\newcommand{\hatl}{\widehat\lambda}
\newcommand{\decl}{\lambda^{(\ge 2)}}
\newcommand{\supp}{\mathrm{supp}}
\newcommand{\SYT}{\mathrm{SYT}}
\newcommand{\Rp}{{R'}}

\title{Extended sign-kill for $\hatl=(4,2)$: \\
       the $E_j^-$ chain at $j\le q-3$ via three-branch support}
\author{Clio}
\date{14 May 2026 (very late prove session, sixth pass)}

\begin{document}
\maketitle

\begin{abstract}
Let $\lambda = (4, 2, 1^q) \vdash n = q+6$ with $q \ge 3$, $\ell = q+2$,
and $B = \Bdec\Vlam = R'_{\ell+1}R'_{\ell+2}R'_{\ell+3}R'_{\ell+4}\,\Vlam$.
We prove the ``$\subseteq$'' direction of the extended sign-kill
conjecture for this family at $q \ge 4$:
\[
   B \;\subseteq\; \bigcap_{j=1}^{q-3} E_j^-\!\mid_{\Vlam},
\]
matching the predicted threshold $\tau(\lambda) = q + 3 -
\hatl_1 - \hatl_2 = q - 3$.  We further prove sharpness at the
boundary $j = q - 2$ (structural reduction + computational
witness for $q \in \{4, 5, 6, 7\}$).  At $q = 3$ the containment
is vacuous (the predicted intersection is empty).

The proof follows the recipe established in the previous (very late
prove session, fifth pass) paper on $\hatl = (3, 2)$.  Three new
ingredients are needed:
\begin{enumerate}[topsep=0.3em,itemsep=0.1em,leftmargin=2em]
\item An extension of the May-20 three-branch reduction to
      $\lambda = (4, 2, 1^q)$, which involves \emph{four} $R'$-factors
      in $\Om$ and hence \emph{three} outer $(T_j{+}1)$ factors after
      $\Phi_X$-extraction (rather than two as in $\hatl = (3, 2)$).
\item Application of the May-14 fifth-pass extended sign-kill
      theorem for $\hatl = (3, 2)$ \emph{recursively} to the inner
      piece $\mu_A = (3, 2, 1^{q-1})$, which requires $q \ge 4$.
\item Application of the hook support lemma (May-14 fifth pass,
      Theorem~3) to the inner piece $\mu_B = (4, 1^q)$, which
      requires $q \ge 4$ ($m \ge k$ at $k = 4$, $m = q$).
\end{enumerate}
The $\mu_C$-piece vanishes by the same column-$1$ leakage mechanism
as in the May-20 paper.  Combined with the universal $B \subseteq
E_\ell^+$ (May-13 evening), this closes the $\hatl = (4, 2)$ row of
the May-14 evening conjecture table in the containment direction at
$q \ge 4$.
\end{abstract}

\section{Setup}\label{sec:setup}

We use the conventions of the May-13--May-14 series.  $\Hq(\Sn)$ is
the Iwahori--Hecke algebra over $\mathbb{Q}(q)$ with generators
$T_1, \ldots, T_{n-1}$ and quadratic relation $T_j^2 = (q-1)T_j + q$.
$\Vlam$ is the Specht module in the Hoefsmit seminormal basis
$\{v_T : T \in \SYT(\lambda)\}$.

For $1 \le j \le n-1$, set $S_j := T_j + 1$.  Write
\[
   E_j^+ \;:=\; \ker(T_j - q)\!\mid_{\Vlam}
   \;=\; \mathrm{Im}(S_j\!\mid_{\Vlam}),
   \qquad
   E_j^- \;:=\; \ker(S_j)\!\mid_{\Vlam}.
\]
For $k \ge 1$, let $R'_k := S_{k-1}\,S_{k-2}\cdots S_1$ with the
convention $R'_1 = 1$.  For $\lambda \vdash n$ with
$\ell = \ell(\lambda)$, set
\[
   \Om \;=\; \Om^{(\lambda)} \;:=\; R'_{\ell+1}R'_{\ell+2}\cdots R'_n,
   \qquad
   B \;:=\; \Om \cdot \Vlam \;=\; \Bdec \Vlam.
\]

For $\lambda = (4, 2, 1^q)$, $n = q + 6$ and $\ell = q + 2$, so
$\Om = R'_{q+3}R'_{q+4}R'_{q+5}R'_{q+6}$ has \emph{four} $R'$-blocks
(equivalently, $\Om = R'_{n-3}R'_{n-2}R'_{n-1}R'_n$).  By the
May-13 multi-corner-dimension paper, $\dim B = f^{\lambda^{(\ge 2)}}
= f^{(3, 1)} = 3$ in the stable regime $q \ge 2$.

\paragraph{Seminormal support.}  For $W \subseteq V_\lambda$, define
\[
   \supp(W) \;:=\; \{\,T \in \SYT(\lambda) : \text{some } w \in W
   \text{ has } [v_T]\,w \ne 0\,\}.
\]

\paragraph{Column-$1$ identity prefix.}  For an SYT $T$ of any
partition $\mu$, the \emph{identity prefix length} is
\[
   p(T) \;:=\; \max\{\,P \ge 0 : T(i, 1) = i \text{ for all }
   1 \le i \le P\,\}.
\]

\section{Input from prior papers}\label{sec:input}

\subsection{The support rule and prefix corollary (May-14 fourth pass)}

\begin{proposition}[Support rule]\label{prop:support-rule}
Let $W \subseteq V_\lambda$ be any subspace.  Fix $1 \le j \le n-1$.
If every $T \in \supp(W)$ has $j$ and $j+1$ in the same column of
$T$, then $W \subseteq E_j^-\!\mid_{V_\lambda}$.
\end{proposition}

\begin{remark}\label{rem:prefix-implies-Ej-}
If every $T \in \supp(W)$ has $p(T) \ge P$, then for every
$1 \le j \le P - 1$ the entries $j, j+1$ sit at column-$1$ rows
$j, j+1$, hence share column~$1$.  By Proposition~\ref{prop:support-rule},
$W \subseteq E_j^-\!\mid_{V_\lambda}$ for $1 \le j \le P - 1$.
\end{remark}

\subsection{The May-14 fifth-pass results}

We need both of the May-14 fifth-pass theorems.

\begin{theorem}[Extended sign-kill for $\hatl = (3, 2)$, May-14 fifth pass]\label{thm:32-prefix}
For $\lambda' = (3, 2, 1^{q'}) \vdash n' = q' + 5$ with $q' \ge 3$,
every $T' \in \supp\bigl(B^{(\lambda')}_{q'+3} V_{\lambda'}\bigr)$
satisfies $p(T') \ge n' - 6 = q' - 1$.  Consequently,
$B^{(\lambda')}_{q'+3} V_{\lambda'} \subseteq E_j^-\!\mid_{V_{\lambda'}}$
for $1 \le j \le q' - 2$.
\end{theorem}

This is Theorem~6 (Main) of \texttt{2026-05-14-extended-sign-kill-32.tex}
together with the prefix-bound construction in its proof of Step~1--Step~3.

\begin{theorem}[Hook support lemma, May-14 fifth pass]\label{thm:hook-support}
Let $\mu = (k, 1^m) \vdash n_0 := k + m$ with $k \ge 2$ and $m \ge k$.
Set $\ell_0 := \ell(\mu) = m + 1$ and $j_0 := \ell_0 + 1 = m + 2$.
Then every $T \in \supp(B^{(\mu)}_{j_0} V_\mu)$ satisfies
$p(T) \ge m - k + 2$.
\end{theorem}

This is Theorem~3 of \texttt{2026-05-14-extended-sign-kill-32.tex}.

\subsection{The outer-factor preservation lemma (May-14 fifth pass)}

\begin{lemma}[Outer-factor preservation]\label{lem:outer-preserve}
Let $W \subseteq V_\nu$ be a subspace such that every $T \in \supp(W)$
satisfies $p(T) \ge P$.  Then for every $j \ge P + 1$,
every $T'' \in \supp((T_j{+}1) W)$ also satisfies $p(T'') \ge P$.
\end{lemma}

This is Lemma~4 of \texttt{2026-05-14-extended-sign-kill-32.tex}.
The proof is elementary: $(T_j+1)$ acts on the seminormal basis
within $\{v_T, v_{s_j T}\}$; the swap $s_j$ moves only entries $j$
and $j+1$, both $\ge j \ge P + 1 > P$, so neither lies in the
column-$1$ prefix and the prefix is preserved.

\subsection{The May-13 hook column-$1$ result (rank-zero)}

We will use the May-12 hook rank-zero theorem:

\begin{theorem}[Hook column-$1$, May-12]\label{thm:hook-E1minus}
For $\mu = (k, 1^m) \vdash n_0$ with $k \ge 2$ and $m \ge \max(k, 2)$,
$B^{(\mu)}_{\ell_0+1} V_\mu \subseteq E_1^-\!\mid_{V_\mu}$.
\end{theorem}

This is the May-12 hook paper (\texttt{2026-05-12-hook-j-formula.tex}),
also restated in the May-20 paper as Theorem~2.3.

\subsection{Phase A endpoint (May-19)}

\begin{theorem}[Abstract Phase A endpoint]\label{thm:phaseA}
For any $\Hq$-module $V$ and any $1 \le j \le n-1$,
$(T_j+1)\cdots(T_1+1) V = E_j^+(V)$.  In particular,
$R'_n V_\lambda = E_{n-1}^+\!\mid_{V_\lambda}$.
\end{theorem}

\section{The three-branch reduction for $(4, 2, 1^q)$}\label{sec:reduction}

For $\lambda = (4, 2, 1^q)$ with $q \ge 3$ (so $n \ge 9$, $\ell = q+2
= n-4$), the three removable corners are
\[
   c_1 = (1, 4), \qquad c_2 = (2, 2), \qquad c_3 = (n-4, 1).
\]
The three doubly-stripped partitions are
\begin{align*}
   \mu_A &:= \lambda \setminus \{c_1, c_3\} = (3, 2, 1^{q-1}), \\
   \mu_B &:= \lambda \setminus \{c_2, c_3\} = (4, 1^q), \\
   \mu_C &:= \lambda \setminus \{c_1, c_2\} = (3, 1^{q+1}),
\end{align*}
each of size $n - 2$.

\paragraph{The branching maps $\Phi_X$.}
Exactly as in the May-20 paper, for each $X \in \{A, B, C\}$ and each
$T' \in \SYT(\mu_X)$, define two SYTs $T_A^X, T_B^X$ of $\lambda$
by extending $T'$ with $\{n-1, n\}$ at the corresponding corner pair
in the two possible orders, and set
\[
   \Phi_X(v_{T'}) \;:=\; v_{T^X_A} \,+\, \tau_X\,v_{T^X_B},
\]
where $\tau_X = \tau(d_X) \in \mathbb{Q}(q)^\times$ is the Hoefsmit
$T_{n-1}$-block coefficient at axial distance
$d_X := c(c_X^+) - c(c_X^-)$.  The axial distances are
\[
   |d_A| \;=\; |(1-(n-4))-(4-1)| \;=\; n - 6, \qquad
   |d_B| \;=\; |(1-(n-4))-(2-2)| \;=\; n - 5, \qquad
   |d_C| \;=\; |0 - 3| \;=\; 3.
\]
All $\ge 3 \ge 2$ for $q \ge 3$, so the $T_{n-1}$ $2$-blocks are
non-degenerate.

\begin{lemma}[$\Phi_X$ properties]\label{lem:Phi-prop}
For each $X \in \{A, B, C\}$:
\begin{enumerate}[topsep=0.3em,itemsep=0.1em,leftmargin=2em,
                  label=(\roman*)]
\item $\Phi_X: V_{\mu_X} \hookrightarrow V_\lambda$ is injective and
   $T_i$-equivariant for $1 \le i \le n - 3$.
\item The seminormal supports of $\Phi_A(V_{\mu_A}),
   \Phi_B(V_{\mu_B}), \Phi_C(V_{\mu_C})$ are pairwise disjoint and
   their direct sum equals $E_{n-1}^+\!\mid_{V_\lambda}$.
\end{enumerate}
\end{lemma}

\begin{proof}
The argument is verbatim May-20 Lemmas~3.2--3.4: injectivity from
distinct corner-pair extensions; equivariance from the
$T_{n-1}$-commutation with $\{T_1, \ldots, T_{n-3}\}$ and the
decoupling argument on $4$-dimensional joint blocks; disjointness
from distinct corner-pair supports; dimension match by counting
$+q$-eigenvectors over the three $T_{n-1}$-corner-pair classes.
The argument depends only on $\lambda$ having three corners and on
the axial distances being $\ge 2$, both of which hold for
$\lambda = (4, 2, 1^q)$ at every $q \ge 3$.
\qed
\end{proof}

\begin{theorem}[Three-branch reduction for $(4, 2, 1^q)$]\label{thm:reduction}
For $\lambda = (4, 2, 1^q)$ with $q \ge 3$,
\begin{equation}\label{eq:reduction}
   B \;=\; (T_{n-4}{+}1)(T_{n-3}{+}1)(T_{n-2}{+}1)\,
   \sum_{X \in \{A,B,C\}}
   \Phi_X\bigl(B_{n-4}^{(\mu_X)} V_{\mu_X}\bigr).
\end{equation}
\end{theorem}

\begin{proof}
We follow the May-15/May-20 pull-through argument, adapted to
$\Om = R'_{n-3}R'_{n-2}R'_{n-1}R'_n$ (four $R'$-factors).

\paragraph{Step 1 (Phase A).}  By Theorem~\ref{thm:phaseA} at
$V = V_\lambda$, $R'_n V_\lambda = E_{n-1}^+\!\mid_{V_\lambda}$.
By Lemma~\ref{lem:Phi-prop}(ii), this equals
$\bigoplus_X \Phi_X(V_{\mu_X})$.

\paragraph{Step 2 (apply $R'_{n-1}$).}
Write $R'_{n-1} = (T_{n-2}{+}1)\,U$ where $U := (T_{n-3}{+}1)\cdots
(T_1{+}1) \in \Hq(S_{n-2})$.  Since $T_1, \ldots, T_{n-3}$ commute
with $T_{n-1}$ (index gaps $\ge 2$), $U$ preserves
$E_{n-1}^+\!\mid_{V_\lambda}$.  By Lemma~\ref{lem:Phi-prop}(i),
$U\,\Phi_X(v) = \Phi_X(\Rp_{n-2}^{(\mu_X)}\,v)$ for $v \in V_{\mu_X}$
(here $\Rp_{n-2}^{(\mu_X)}$ is the analogue of $R'_{n-2}$ inside
$\Hq(S_{n-2})$, using generators $T_1, \ldots, T_{n-3}$).  Hence
\[
   R'_{n-1} R'_n V_\lambda
   \;=\; (T_{n-2}{+}1)\,\sum_X
        \Phi_X\bigl(\Rp_{n-2}^{(\mu_X)} V_{\mu_X}\bigr).
\]

\paragraph{Step 3 (apply $R'_{n-2}$).}  Write $R'_{n-2} =
(T_{n-3}{+}1)\,S$ where $S := (T_{n-4}{+}1)\cdots(T_1{+}1) \in
\Hq(S_{n-3})$.  $S$ commutes with $T_{n-2}$ (index gap $\ge 2$), so
$R'_{n-2}(T_{n-2}{+}1) = (T_{n-3}{+}1)(T_{n-2}{+}1)\,S$.  And $S
\in \Hq(S_{n-3}) \subset \Hq(S_{n-2})$ acts on $\Phi_X(v)$ as
$\Phi_X(\Rp_{n-3}^{(\mu_X)}\,v)$.  Therefore
\[
   R'_{n-2} R'_{n-1} R'_n V_\lambda
   \;=\; (T_{n-3}{+}1)(T_{n-2}{+}1)\,\sum_X
        \Phi_X\bigl(\Rp_{n-3}^{(\mu_X)} \Rp_{n-2}^{(\mu_X)} V_{\mu_X}\bigr).
\]

\paragraph{Step 4 (apply $R'_{n-3}$).}  Write $R'_{n-3} =
(T_{n-4}{+}1)\,S'$ where $S' := (T_{n-5}{+}1)\cdots(T_1{+}1) \in
\Hq(S_{n-4})$.  $S'$ commutes with both $T_{n-3}$ (gap $= 2$) and
$T_{n-2}$ (gap $= 3$), so
$R'_{n-3}(T_{n-3}{+}1)(T_{n-2}{+}1) = (T_{n-4}{+}1)(T_{n-3}{+}1)(T_{n-2}{+}1)\,S'$.
$S' \in \Hq(S_{n-4}) \subset \Hq(S_{n-2})$ acts on $\Phi_X(v)$ as
$\Phi_X(\Rp_{n-4}^{(\mu_X)}\,v)$.  Combining:
\[
   B \;=\; \Om V_\lambda
   \;=\; (T_{n-4}{+}1)(T_{n-3}{+}1)(T_{n-2}{+}1)\,\sum_X
        \Phi_X\bigl(\Rp_{n-4}^{(\mu_X)} \Rp_{n-3}^{(\mu_X)} \Rp_{n-2}^{(\mu_X)} V_{\mu_X}\bigr).
\]
The triple product inside is exactly $B_{n-4}^{(\mu_X)} V_{\mu_X}$
(three consecutive $R'$-factors at indices $n-4, n-3, n-2$ in
$V_{\mu_X}$, which has $n - 2$ letters), and~\eqref{eq:reduction}
follows.
\qed
\end{proof}

\section{The $C$-piece vanishes}\label{sec:C-zero}

\begin{lemma}[$C$-piece vanishing]\label{lem:C-zero}
$B_{n-4}^{(\mu_C)} V_{\mu_C} = 0$ for $\lambda = (4, 2, 1^q)$ with
$q \ge 3$.
\end{lemma}

\begin{proof}
$\mu_C = (3, 1^{q+1})$ has length $\ell(\mu_C) = q + 2 = n - 4$ and
size $n - 2$, so $j_0(\mu_C) = \ell(\mu_C) + 1 = n - 3$.  By
Theorem~\ref{thm:hook-E1minus} (applied at $k = 3$, $m = q + 1 \ge
4$, so $m \ge k = 3$),
$B^{(\mu_C)}_{n-3} V_{\mu_C} \subseteq E_1^-\!\mid_{V_{\mu_C}}$.

Now
\[
   B^{(\mu_C)}_{n-4} V_{\mu_C}
   \;=\; \Rp^{(\mu_C)}_{n-4}\,B^{(\mu_C)}_{n-3} V_{\mu_C}.
\]
The rightmost factor of $\Rp^{(\mu_C)}_{n-4}$ is $(T_1{+}1)$, which
annihilates $E_1^-\!\mid_{V_{\mu_C}}$.  Hence
$B^{(\mu_C)}_{n-4} V_{\mu_C} = 0$.
\qed
\end{proof}

\begin{remark}\label{rem:C-mechanism}
As in the May-20 paper for $(3, 2, 1^q)$, the $C$-corner pair
($\{c_1, c_2\}$) avoids the column-$1$ bottom corner $c_3$.  This
strips two cells from $\lambda$ without consuming the column-$1$
bottom, producing a hook $\mu_C$ that is one row \emph{shorter}
than $\mu_A, \mu_B$ (i.e., $\ell(\mu_C) = q + 2$ vs.\
$\ell(\mu_A) = \ell(\mu_B) = q + 1$).  The shorter $\mu_C$ has its
$j$-formula at $j_0(\mu_C) = n - 3$ rather than $n - 4$; the index
$n - 4$ at which we evaluate sits one step \emph{below}, and the
leading $(T_1{+}1)$ of $\Rp^{(\mu_C)}_{n-4}$ kills the
$E_1^-$-contained result.  This is the same structural mechanism
as in $\hatl = (3, 2)$.
\end{remark}

By Lemma~\ref{lem:C-zero}, the reduction~\eqref{eq:reduction}
simplifies for $\lambda = (4, 2, 1^q)$, $q \ge 3$, to
\begin{equation}\label{eq:reduction-AB}
   B \;=\; (T_{n-4}{+}1)(T_{n-3}{+}1)(T_{n-2}{+}1)\,
   \Bigl(\Phi_A\bigl(B_{n-4}^{(\mu_A)} V_{\mu_A}\bigr)
   + \Phi_B\bigl(B_{n-4}^{(\mu_B)} V_{\mu_B}\bigr)\Bigr).
\end{equation}

\section{Main theorem}\label{sec:main}

\begin{theorem}[Main]\label{thm:main}
For $\lambda = (4, 2, 1^q) \vdash n = q + 6$ with $q \ge 4$,
\[
   B \;\subseteq\; \bigcap_{j=1}^{q-3} E_j^-\!\mid_{V_\lambda}.
\]
\end{theorem}

\begin{proof}
By Remark~\ref{rem:prefix-implies-Ej-}, it suffices to show that
every $T \in \supp(B)$ satisfies $p(T) \ge q - 2 = n - 8$.  We use
the reduction~\eqref{eq:reduction-AB} together with prefix bounds
on the two inner pieces and Lemma~\ref{lem:outer-preserve} on the
three outer factors.

\medskip\textbf{Step 1: prefix bound on the $\Phi_A$-image.}
$\mu_A = (3, 2, 1^{q-1})$ has size $n - 2 = q + 4$, so $n' = q + 4$
and $q' = q - 1$.  Since $q \ge 4$ we have $q' \ge 3$, so
Theorem~\ref{thm:32-prefix} applies and gives
\[
   p(T') \;\ge\; n' - 6 \;=\; q - 2 \;=\; n - 8
   \qquad \text{for every } T' \in \supp(B^{(\mu_A)}_{n-4} V_{\mu_A}).
\]
(Here $B^{(\mu_A)}_{n-4} = B^{(\mu_A)}_{\ell(\mu_A) + 1}$ since
$\ell(\mu_A) = q + 1 = n - 5$.)

Now $\Phi_A(v_{T'})$ is supported on two SYTs $T_A^A, T_B^A$ of
$\lambda$ obtained from $T'$ by placing $\{n-1, n\}$ at $\{c_1, c_3\}
= \{(1, 4), (n-4, 1)\}$ in the two possible orders.  Cells of $\mu_A$'s
column~$1$ are at rows $1, 2, \ldots, \ell(\mu_A) = n - 5$.  The
added cell $(1, 4)$ is in row~$1$ but column~$4$ (not column~$1$),
and $(n-4, 1)$ is at row $n - 4 > n - 5$, i.e., \emph{below} $\mu_A$'s
column-$1$ cells.  Hence rows $1, \ldots, n - 5$ of column~$1$ in
$T_A^A, T_B^A$ are unchanged from $T'$, so
$p(T_A^A), p(T_B^A) \ge n - 8$.

\medskip\textbf{Step 2: prefix bound on the $\Phi_B$-image.}
$\mu_B = (4, 1^q)$, a hook with $k = 4$, $m = q$.  Since $q \ge 4$
we have $m \ge k$, so Theorem~\ref{thm:hook-support} applies and
gives
\[
   p(T') \;\ge\; m - k + 2 \;=\; q - 2 \;=\; n - 8
   \qquad \text{for every } T' \in \supp(B^{(\mu_B)}_{n-4} V_{\mu_B}).
\]
(Again $B^{(\mu_B)}_{n-4} = B^{(\mu_B)}_{\ell(\mu_B)+1}$ since
$\ell(\mu_B) = q + 1 = n - 5$.)

Now $\Phi_B(v_{T'})$ is supported on two SYTs $T_A^B, T_B^B$ obtained
from $T'$ by placing $\{n-1, n\}$ at $\{c_2, c_3\} = \{(2, 2),
(n-4, 1)\}$.  The added cells are at column~$2$ and row $n - 4 >
n - 5 = \ell(\mu_B)$, so they leave the column-$1$ prefix of $T'$
intact.  Hence $p(T_A^B), p(T_B^B) \ge n - 8$.

\medskip\textbf{Step 3: outer factors preserve the prefix.}
The three outer factors in~\eqref{eq:reduction-AB} are $(T_j{+}1)$
for $j \in \{n-4, n-3, n-2\}$.  Each index satisfies $j \ge n - 4 =
(n - 8) + 4 > (n - 8) + 1$, so the hypothesis $j \ge P + 1$ of
Lemma~\ref{lem:outer-preserve} (with $P = n - 8$) is met.  Applying
the lemma three times in sequence, the prefix bound $\ge n - 8$ is
preserved.

\medskip\textbf{Conclusion.}  Combining Steps~1--3 and~\eqref{eq:reduction-AB}:
every $T \in \supp(B)$ has $p(T) \ge n - 8 = q - 2$.  By
Remark~\ref{rem:prefix-implies-Ej-}, $B \subseteq E_j^-$ for $1 \le
j \le q - 3$.
\qed
\end{proof}

\begin{remark}[$q = 3$]\label{rem:q3-vacuous}
At $q = 3$, the predicted threshold is $\tau = q - 3 = 0$ and the
intersection $\bigcap_{j=1}^0 E_j^-$ is the whole module $V_\lambda$,
so the containment is vacuous.  Indeed Theorem~\ref{thm:main}'s
hypothesis $q \ge 4$ is exactly the joint requirement of the
recursive $(3, 2)$ input ($q' \ge 3$, i.e., $q \ge 4$) and the hook
lemma ($m \ge k$, i.e., $q \ge 4$).  At $q = 3$, neither inner
ingredient applies in its stated form, but the prefix bound
$p \ge n - 8 = 1$ holds trivially (every SYT has prefix $\ge 1$),
so the proof framework collapses to a vacuous statement, consistent
with the predicted $\tau = 0$.
\end{remark}

\section{Sharpness}\label{sec:sharp}

\begin{lemma}[Sharpness reduction]\label{lem:sharp-reduction}
Suppose $B \subseteq E_{q-2}^-\!\mid_{V_\lambda}$, where $\lambda =
(4, 2, 1^q)$ with $q \ge 4$.  Then every $T \in \supp(B)$ satisfies
$p(T) \ge q - 1 = n - 7$.
\end{lemma}

\begin{proof}
Set $j_0 := q - 2 = n - 8$.  Suppose for contradiction $T \in
\supp(B)$ has $p(T) = n - 8$ exactly.  Since $T \in \supp(B)$,
there exists $w \in B$ with $C_T := [v_T]\,w \ne 0$; by hypothesis,
$T_{j_0}\,w = -w$.

\medskip\textbf{Locating entry $n - 7$ in $T$.}  By $p(T) \ge n - 8$,
the cells $\{(1, 1), \ldots, (n - 8, 1)\}$ are filled with entries
$\{1, \ldots, n - 8\}$.  Since $p(T) = n - 8$ exactly,
$T(n - 7, 1) \ne n - 7$.  The remaining $8$ cells of $\lambda$ are
\[
   \{(1, 2),(1, 3),(1, 4),(2, 2),(n-7,1),(n-6,1),(n-5,1),(n-4,1)\},
\]
holding the $8$ entries $\{n-7, n-6, \ldots, n\}$.

We claim entry $n - 7$ is forced to sit at $(1, 2)$:
\begin{itemize}[topsep=0.3em,itemsep=0.1em]
\item Not at $(n - 6, 1), (n - 5, 1), (n - 4, 1)$: column-$1$
   increase forces $T(n - 7, 1) < n - 7$, but no eligible entry
   $< n - 7$ remains for $(n - 7, 1)$.
\item Not at $(1, 3)$: $T(1, 2) < T(1, 3) = n - 7$, so $T(1, 2) \in
   \{n - 7\} \cap (\text{remaining})$ — only candidate is $T(1, 2) <
   n - 7$, but all remaining entries are $\ge n - 7$.  Impossible.
\item Not at $(1, 4)$: same argument, $T(1, 2), T(1, 3) < n - 7$
   forced, both impossible.
\item Not at $(2, 2)$: column-$2$ increase forces $T(1, 2) < n - 7$,
   impossible since $T(1, 2) \ge n - 7$ over the remaining entries.
\end{itemize}
Hence $T(1, 2) = n - 7$.

\medskip\textbf{Off-diagonal swap.}  The cells of $j_0 = n - 8$ (at
$(n - 8, 1)$) and $j_0 + 1 = n - 7$ (at $(1, 2)$) are in different
rows ($n - 8 \ne 1$ for $n \ge 10$, i.e., $q \ge 4$) and different
columns.  Axial distance:
\[
   |d| \;=\; |\,c(1, 2) - c(n-8, 1)\,|
   \;=\; |\,1 - (1 - (n-8))\,|
   \;=\; n - 8 \;\ge\; 2 \quad \text{for } q \ge 4.
\]
Hoefsmit's $2$-block at distance $|d| \ge 2$ is non-degenerate:
$T_{j_0}\,v_T = \alpha\,v_T + \beta\,v_{T^*}$ with $\beta \ne 0$,
where $T^* := s_{j_0}T$ swaps entries $n - 8$ and $n - 7$ between
cells $(n - 8, 1)$ and $(1, 2)$.

\medskip\textbf{$T^*$ is a valid SYT.}  After the swap, $(n-8, 1)$
holds $n - 7$ and $(1, 2)$ holds $n - 8$.  Column-$1$ rows
$1, \ldots, n - 9$ unchanged ($1, \ldots, n - 9$); $(n - 8, 1) =
n - 7 > T(n - 9, 1) = n - 9$ \checkmark; $T(n - 7, 1) > T(n - 8, 1)|_{T^*}
= n - 7$ \checkmark (we have $T(n - 7, 1) \in \{n - 6, \ldots, n\}$ from the
forced argument above).  Row~$1$: $T(1, 1) = 1 < T(1, 2)|_{T^*} =
n - 8$ \checkmark; $T(1, 2)|_{T^*} = n - 8 < T(1, 3) \in \{n - 6, \ldots, n\}$
\checkmark.  Column~$2$: $T(1, 2)|_{T^*} = n - 8 < T(2, 2) \in
\{n - 6, \ldots, n\}$ \checkmark.  So $T^*$ is a valid SYT of $\lambda$.

\medskip\textbf{Prefix drop forces $T^* \notin \supp(B)$.}
$T^*(n - 8, 1) = n - 7 \ne n - 8$, so $p(T^*) \le n - 9 < n - 8$.
By Theorem~\ref{thm:main}, every $U \in \supp(B)$ has
$p(U) \ge n - 8$, hence $T^* \notin \supp(B) \supseteq \supp(w)$.

\medskip\textbf{Coefficient contradiction.}  We compute
\[
   [v_{T^*}]\,(T_{j_0}\,w)
   \;=\; \sum_{U \in \supp(w)} C_U \cdot [v_{T^*}]\,(T_{j_0}\,v_U).
\]
The diagonal contribution $[v_{T^*}]\,(T_{j_0}\,v_{T^*}) =$ (an
$\alpha$-multiple of $C_{T^*}$) requires $T^* \in \supp(w)$ (no).
The off-diagonal contribution at $v_U$ requires the $T_{j_0}$
$2$-block partner of $U$ to be $T^*$, i.e., $s_{j_0} U = T^*$,
i.e., $U = s_{j_0} T^* = T$.  Hence
\[
   [v_{T^*}]\,(T_{j_0}\,w) \;=\; C_T \cdot \beta \;\ne\; 0,
\]
while $[v_{T^*}]\,(-w) = -C_{T^*} = 0$.  This contradicts
$T_{j_0}\,w = -w$.  We conclude no $T \in \supp(B)$ has
$p(T) = n - 8$ exactly, hence every such $T$ has $p(T) \ge n - 7$.
\qed
\end{proof}

\begin{theorem}[Sharpness, computationally verified]\label{thm:sharp}
For $\lambda = (4, 2, 1^q)$ at every $q \in \{4, 5, 6, 7\}$,
$\supp(B)$ contains an SYT $T$ with $p(T) = n - 8 = q - 2$ exactly;
consequently $B \not\subseteq E_{q-2}^-\!\mid_{V_\lambda}$.
\end{theorem}

\begin{proof}
Existence of $T \in \supp(B)$ with $p(T) = q - 2$ is verified
computationally for $q \in \{4, 5, 6, 7\}$:
\texttt{find\_witness.py} (Section~\ref{sec:verify}) reports
exactly $42$ such witnesses in each case.  An explicit witness at
$q = 4$ ($n = 10$, prefix $= 2$):
\[
   T(1, \cdot) = (1, 3, 4, 5), \quad T(2, \cdot) = (2, 6), \quad
   T(3, 1) = 7,\ T(4, 1) = 8,\ T(5, 1) = 9,\ T(6, 1) = 10.
\]
Column~$1$ entries are $(1, 2, 7, 8, 9, 10)$, so $p(T) = 2$ exactly.
(Analogous witnesses for $q = 5, 6, 7$ are listed by
\texttt{find\_witness.py}.)  Combining with
Lemma~\ref{lem:sharp-reduction}, the assumption $B \subseteq
E_{q-2}^-$ would force every $U \in \supp(B)$ to have
$p(U) \ge q - 1$, contradicting the witness $T$ with
$p(T) = q - 2$.  Hence $B \not\subseteq E_{q-2}^-$.
\qed
\end{proof}

\begin{remark}[Stronger empirical statement]\label{rem:stronger-sharp}
Computational verification (mod $P = 100\,003$, $Q = 11$) at
$q \in \{4, 5, 6, 7\}$ shows that $S_{j_0}\!\mid_B$ has full rank
$3 = \dim B$ at $j_0 = q - 2$, i.e., $B \cap E_{j_0}^- = 0$ (not
merely $B \not\subseteq E_{j_0}^-$).  Theorem~\ref{thm:sharp} proves
the weaker non-containment.
\end{remark}

\section{Eigenspace summary}\label{sec:summary}

\begin{theorem}[Eigenspace summary, $\hatl = (4, 2)$]\label{thm:summary}
For $\lambda = (4, 2, 1^q)$ with $q \ge 4$,
\[
   B \;\subseteq\; E_\ell^+\;\cap\;\bigcap_{j=1}^{q-3} E_j^-\!\mid_{V_\lambda},
   \qquad
   B \not\subseteq E_{q-2}^-\!\mid_{V_\lambda}.
\]
The $E_\ell^+$ containment is sharp with scalar action: $T_\ell\!\mid_B
= q\!\cdot\!\mathrm{Id}_B$ (May-13 evening Theorem~A).
\end{theorem}

\begin{proof}
$E_j^-$ chain at $j \le q - 3$: Theorem~\ref{thm:main}.
Non-containment at $j = q - 2$: Theorem~\ref{thm:sharp}.
$E_\ell^+$ containment with scalar action: May-13 evening (universal
across all $\lambda$ in stable regime).
\qed
\end{proof}

The eigenspace constraints alone do not pin down $B$: computationally
the intersection on the right has dimension strictly $> \dim B = 3$
in every tested case.  The dimension formula
$\dim B = f^{\lambda^{(\ge 2)}} = f^{(3, 1)} = 3$ (May-13
multi-corner) is a sharper fact.

\section{Computational verification}\label{sec:verify}

\begin{theorem}[Verification]\label{thm:verify}
Theorem~\ref{thm:summary} was verified computationally at $q \in
\{2, 3, 4, 5, 6, 7\}$ (i.e., $n \in \{8, 9, 10, 11, 12, 13\}$),
mod $P = 100\,003$ at $Q = 11$.  In every case:
\begin{itemize}[topsep=0.3em,itemsep=0.1em]
\item $\dim B = 3$, matching the May-13 multi-corner formula.
\item $|\supp(B)|$ stabilizes at $123$ for $\lambda = (4, 2, 1^q)$
   at $q \ge 4$ (boundary value $|\supp(B)| = 115$ at $q = 3$,
   $69$ at $q = 2$).  Cf.\ $|\supp(B)| = 26$ for $(3, 2, 1^q)$ in
   the May-14 fifth-pass paper.
\item Every $T \in \supp(B)$ has $p(T) \ge q - 2 = n - 8$, with the
   bound achieved at $42$ SYTs per $q$ (the sharpness witnesses,
   for $q \in \{4, 5, 6, 7\}$).
\item $T_j\,b = -b$ for every $b \in B$ and $1 \le j \le q - 3$.
\item $S_{q-2}\!\mid_B$ has rank $\dim B = 3$ (stronger than
   Theorem~\ref{thm:sharp}; cf.\ Remark~\ref{rem:stronger-sharp}).
\item The $C$-piece $B_{n-4}^{(\mu_C)} V_{\mu_C} = 0$ verified
   at every $q \ge 3$.
\item Inner pieces: $\dim B^{(\mu_A)} = 2$, $\dim B^{(\mu_B)} = 1$
   (sum $= 3 = \dim B$, matching dim-additivity), with min prefix
   $= q - 2$ for both.
\end{itemize}
\end{theorem}

\begin{proof}
Scripts in
\texttt{\textasciitilde/projects/scratch/2026-05-14-extended-sign-kill-42/}.
The script \texttt{explore\_42\_family.py} constructs Hoefsmit
seminormal matrices for $\lambda$ and each $\mu_X$, builds $\Om$
and $\Om^{(\mu_X)}$, computes supports by row-echelon reduction, and
checks $S_j$-rank by direct linear algebra.  The script
\texttt{find\_witness.py} enumerates SYTs in $\supp(B)$ with
$p(T) = q - 2$ exactly and prints an explicit witness for each $q$.
\qed
\end{proof}

\section{Discussion}\label{sec:discussion}

\subsection{What is new}

This paper closes the $\hatl = (4, 2)$ row of the extended sign-kill
conjecture (May-14 evening) in the containment direction at the
predicted threshold $\tau = q - 3$, and proves sharpness at the
boundary $j = q - 2$ (modulo the computational input
Theorem~\ref{thm:sharp}, verified for $q \in \{4, 5, 6, 7\}$).
Prior to this paper:
\begin{itemize}[topsep=0.3em,itemsep=0.1em]
\item May-14 fourth-pass: $\hatl = (2, 2)$, all $j$.
\item May-14 fifth-pass: $\hatl = (3, 2)$, all $j$ (containment + sharpness).
\end{itemize}

\subsection{What composes here that didn't before}

This is the first case where the recursive ingredient (the May-14
fifth-pass $(3, 2)$ theorem) is invoked as an \emph{input}, applied
to the inner $\mu_A$ piece of a strictly larger shape.  This
validates the compositional strategy implicit in the
``$r$-additivity meta-theorem'' (May-13 evening): the $(4, 2)$
proof inherits its prefix bound from $\mu_A = (3, 2, 1^{q-1})$
(today) and $\mu_B = (4, 1^q)$ (hook lemma), each of which inherits
from its own ingredients.  The May-14 fifth-pass paper closed the
$(3, 2)$ row knowing $(2, 2)$ as input; today's paper closes the
$(4, 2)$ row knowing $(3, 2)$ as input.  This is a clean inductive
ladder.

\subsection{The hypothesis $q \ge 4$}

Both inner ingredients require $q \ge 4$:
\begin{itemize}[topsep=0.3em,itemsep=0.1em]
\item Recursive $(3, 2)$ at $\mu_A = (3, 2, 1^{q-1})$ requires
   $q' = q - 1 \ge 3$, i.e., $q \ge 4$.
\item Hook lemma at $\mu_B = (4, 1^q)$ requires $m \ge k$, i.e.,
   $q \ge 4$.
\end{itemize}
For $q = 3$ the predicted threshold $\tau = 0$ makes the containment
vacuous; for $q = 2$ ($\tau = -1$) likewise.  We verified $\dim B
= 3$ at $q \in \{2, 3\}$ as well, confirming the stable regime
starts no later than $q = 2$.  ($q = 1$ is non-stable, $\dim B = 6
= 2 \cdot 3$ — the factor-of-$2$ row-repeat phenomenon also
investigated and refuted as a uniform conjecture in the May-14
morning paper.)

\subsection{Toward general $\hatl$}

The natural next family is $\hatl = (3, 3)$, i.e., $\lambda = (3,
3, 1^q)$.  Here $\hatl_1 = \hatl_2 = 3$, so the predicted threshold
is $\tau = q + 3 - 3 - 3 = q - 3$, same as $\hatl = (4, 2)$.  The
corner structure differs:
\begin{itemize}[topsep=0.3em,itemsep=0.1em]
\item For $\lambda = (3, 3, 1^q)$, $\lambda_1 = \lambda_2 = 3$, so
   $(1, 3)$ is \emph{not} a removable corner (row 1 = row 2).  The
   removable corners are $c_1 = (2, 3)$ (the bottom of row 2) and
   $c_2 = (q+2, 1)$ (column-$1$ bottom).  Only \emph{two} corners,
   so this is an $r = 1$ shape, with $\dim B = f^{\lambda^{(\ge 2)}}
   = f^{(2, 2)} = 2$.
\end{itemize}
A two-corner shape has a simpler reduction: $r = 1$ in the
multi-corner classification, requiring only the May-15 \emph{single}
branching reduction (one $\Phi$, one inner piece).  The prediction
$\tau = q - 3$ for $\hatl = (3, 3)$ should follow from the
single-branch version of today's framework, with inner piece
$\mu = \lambda \setminus \{c_1, c_2\} = (3, 2, 1^{q-1})$
(recursive!) or possibly $(2, 2, 1^{q-1})$ depending on the corner
labels.  This is structurally simpler than today's three-branch
case and should follow from a single invocation of the $(3, 2)$
recursive theorem.

The next interesting $r = 2$ family is $\hatl = (4, 3)$, i.e.,
$\lambda = (4, 3, 1^q)$.  Predicted $\tau = q - 4$.  Three corners
again ($c_1 = (1, 4)$, $c_2 = (2, 3)$, $c_3 = (q+2, 1)$), three
doubly-stripped:
\begin{itemize}[topsep=0.3em,itemsep=0.1em]
\item $\mu_A = (3, 3, 1^{q-1})$: a $\hatl = (3, 3)$ shape
   (recursive — requires the $\hatl = (3, 3)$ result above).
\item $\mu_B = (4, 2, 1^{q-1})$: a $\hatl = (4, 2)$ shape
   (recursive — today's result at $q - 1$).
\item $\mu_C = (3, 2, 1^q)$: a $\hatl = (3, 2)$ shape (recursive
   — May-14 fifth pass).
\end{itemize}
None of these is a hook!  The hook lemma is not needed for $(4, 3)$;
all three branches reduce to non-hook inner pieces, each handled by
the recursive extended sign-kill chain.  This is the first case
where the proof framework is \emph{fully recursive} (no hook input).

\subsection{What the iso programme this is not}

This paper continues the structural sharpening of $B$ via its
$T_j$-eigenspace containments — a programme orthogonal to (and
unaffected by) the iso programme refutations of May-13 and
May-14 (morning/afternoon).  We learn more about \emph{where} $B$
sits inside $V_\lambda$ via these eigenspace conditions, without
claiming $B$ equals any standard representation-theoretic object.

\section{Status}\label{sec:status}

\paragraph{Established.}
\begin{itemize}[topsep=0.3em,itemsep=0.1em]
\item Theorem~\ref{thm:reduction}: three-branch reduction for
   $\lambda = (4, 2, 1^q)$, $q \ge 3$.
\item Lemma~\ref{lem:C-zero}: $\mu_C$-piece vanishes.
\item Theorem~\ref{thm:main}: $B \subseteq E_j^-$ for $j \le q - 3$
   at $q \ge 4$.
\item Theorem~\ref{thm:sharp}: sharpness at $j = q - 2$ for
   $q \in \{4, 5, 6, 7\}$ (and Lemma~\ref{lem:sharp-reduction}'s
   structural reduction at all $q \ge 4$).
\item Theorem~\ref{thm:summary}: complete eigenspace summary.
\end{itemize}

\paragraph{Open.}
\begin{itemize}[topsep=0.3em,itemsep=0.1em]
\item Structural (non-computational) sharpness across all $q \ge 4$
   (same obstacle as the May-14 fifth-pass paper; the residual is
   existence of a boundary-prefix witness in $\supp(B)$).
\item Sharpness at $j \in \{q - 1, q, \ldots, n - 1\} \setminus
   \{\ell\}$ (consistent computationally; argument should be
   analogous to the May-14 fifth-pass paper).
\item The next family $\hatl = (3, 3)$ at $\lambda = (3, 3, 1^q)$
   (single-branch reduction, recursive on $\mu = (3, 2, 1^{q-1})$).
\item $\hatl = (4, 3)$ at $\lambda = (4, 3, 1^q)$ (fully recursive
   three-branch, no hook input).
\item General $\hatl$, $\ell(\hatl) \ge 2$: the support prefix bound
   $p(T) \ge q + 3 - \hatl_1 - \hatl_2$ across $\supp(B)$
   (May-14 fifth-pass Conjecture~9).
\end{itemize}

\paragraph{Push status.}
Pending PAT refresh; prior unpushed-commit count on
\texttt{clio-vega/proofs}: $66$ (as of May-14 fifth-pass session).

\paragraph{Scripts.}
\begin{itemize}[topsep=0.3em,itemsep=0.1em]
\item \texttt{\textasciitilde/projects/scratch/2026-05-14-extended-sign-kill-42/explore\_42\_family.py}
   --- main verification.
\item \texttt{\textasciitilde/projects/scratch/2026-05-14-extended-sign-kill-42/find\_witness.py}
   --- sharpness witnesses.
\end{itemize}

\end{document}
